\section{Discussion}

\subsection{Interpretation of Results}

The 90\% experimental success rate demonstrates that testability is a predictable property via formal constraint-satisfiability verification in proof-of-concept settings. By validating predictions against post-hoc p-values rather than circular expert agreement, we demonstrate non-circular ground truth for meta-research tool evaluation. The system outperforms random baseline (50\%) with high statistical significance (binomial $p = 0.0002$) and approaches expert judgment levels (80-85\% inter-rater reliability) while avoiding the circularity of expert-consensus validation.

Cross-domain confound pattern transfer (93.33\% precision) reveals that documented confounds reflect fundamental deep learning design choices --- preprocessing-metric coupling, architecture-data coupling, hyperparameter interdependencies --- that generalize beyond modality boundaries. NLP confound patterns (tokenizer-BLEU from Salesky 2020) successfully flag vision confounds (resolution-architecture from Touvron 2019) and training confounds (batch-LR from Goyal 2017), suggesting confound detection can leverage a shared pattern library rather than requiring domain-specific rule engineering for each new modality.

The two null results (testable hypotheses yielding $p = 0.679$ and $p = 0.757$) validate the system's distinction between testability and guaranteed significance. Both hypotheses had valid (D,B,M) triples and no flagged confounds, yet interventions did not produce significant effects. This outcome confirms the conservative design: the system predicts experimental feasibility (resources available, no confounds present) but does not guarantee effect sizes or statistical significance.

\subsection{Limitations}

\textbf{KB Coverage Ceiling (84\%):} Eight well-known datasets are missing from the January 2026 HuggingFace snapshot, causing h-m2 false negatives where testable hypotheses referencing Pascal VOC, STL-10, SST-2, Common Voice, tieredImageNet, CUB-200, or Stanford Cars are incorrectly classified as ``not testable.'' This limitation reflects catalog incompleteness rather than extraction failures. Multi-source KB aggregation (HuggingFace + Papers With Code + TensorFlow Datasets + Google Dataset Search) is projected to achieve 95\%+ coverage through complementary catalog strengths, reducing false negatives from 6/10 to 1-2/10.

\textbf{Confound Pattern Incompleteness (15 patterns, 2017-2020):} The pattern database misses post-2020 confounds such as pre-training dataset confounds (transfer learning), contrastive loss confounds (self-supervised learning), and prompt template confounds (large language models). One h-m3 false negative (``Compare pre-training datasets (same architecture)'') stems from this limitation --- the pattern database lacks pre-training confound patterns from post-2020 transfer learning literature. Living confound database with automated literature mining (arXiv/NeurIPS abstracts 2020-2026) and crowdsourced contributions is proposed to expand coverage.

\textbf{Keyword Extraction Brittleness (10\% recall):} Simple substring matching fails on paraphrased hypotheses where dataset/benchmark/metric names are expressed non-standardly (``improve COCO numbers'' vs. ``improve COCO mAP''). This brittleness causes h-m2 false negatives where testable hypotheses with informal phrasing are missed. Semantic (D,B,M) extraction via sentence-BERT embeddings is projected to improve recall to 70-80\% while maintaining false positive rate below 10\% through tuned similarity thresholds (0.8-0.9).

\textbf{PoC Validation (Simplified Experiments):} h-m4 validation used mock data with known effect sizes to test verification pipeline logic (internal validity) rather than real-world experiments with actual datasets and full training runs (external validity). The 90\% success rate may overestimate real-world performance --- anticipated degradation to 70-80\% due to dataset loading failures, training instabilities, and hyperparameter sensitivity in production settings. Real-world experiment execution (Weights \& Biases API integration, cloud GPU compute, actual dataset loading) is proposed for external validity testing.

\textbf{Usability Untested:} The P2 prediction (80\% of non-expert users can add new (D,B,M) triples in under 10 minutes) was not validated --- no user studies were conducted. The system is validated for automated use (h-m1: 84\% coverage via API extraction without manual KB updates) but extensibility claim for human-in-loop KB expansion is unsubstantiated. User study with 10 participants, standardized task (add triple for given dataset), and timed correctness measurement is proposed for P2 validation.

\subsection{Broader Impact}

\textbf{Positive:} The system reduces wasted research effort in constraint-driven environments (academic labs with limited resources, independent ML practitioners, reproducibility initiatives) by enabling upfront testability verification in minutes rather than post-hoc discovery after weeks of hypothesis formulation. Automated KB construction (84\% coverage) eliminates manual curation bottlenecks, allowing the system to scale with catalog growth.

\textbf{Negative:} Over-reliance on the system may discourage exploratory hypotheses outside catalog coverage (innovation risk). Researchers might dismiss valid novel ideas if the system classifies them as ``not testable'' due to KB gaps rather than fundamental infeasibility. Mitigation: expose partial confidence scores (e.g., 0.7 = borderline, semantic match but not exact), flag borderline cases for manual expert review, provide interpretable explanations (``Dataset not in catalog'' vs. ``Metric requires human evaluation'').

\textbf{Ethical Considerations:} The system operates on meta-research workflows (hypothesis testability classification) with no direct societal harm pathways. Dual-use concern: the system could be misused to dismiss valid research proposals if applied outside its intended constraint-driven context (existing datasets/benchmarks only). Reviewers might inappropriately reject proposals as ``untestable'' without recognizing the proposals introduce new benchmarks or data collection efforts, which fall outside the system's design scope.

\subsection{Future Work}

Direction 1: Multi-source KB aggregation for 95\%+ coverage. Direction 2: Living confound database with automated literature mining and crowdsourced contributions. Direction 3: Semantic (D,B,M) extraction via sentence-BERT for 70-80\% recall at $<$10\% FPR. Direction 4: Real-world experiment execution for external validity validation (anticipated: 90\% PoC $\rightarrow$ 70-80\% real-world). Direction 5: User study for KB extensibility validation (P2 claim). Direction 6: Partial confidence scoring to expose uncertainty on borderline cases.

Beyond incremental improvements, the vision is deployment-ready testability verification integrated into research workflow tools (hypothesis management systems, experiment tracking platforms, research proposal review systems). The ultimate goal: reduce the rate at which hypotheses are discovered to be infeasible post-hoc through systematic upfront constraint checking, enabling resource-constrained researchers to allocate time efficiently toward experimentally feasible hypotheses.
